\section{Introduction}
\label{sec:intro}
Large language models (LLMs) have been shown deficient in \textit{implicit} reasoning with their parametric memory of knowledge and rules. For example, a range of LLMs are found to be incapable of robustly composing internalized facts~\citep{press-etal-2023-measuring, yang2024large}, and even GPT-4~\citep{openai2023gpt4} cannot adequately compare entities' attributes despite knowing them~\citep{allenzhu2023physics}.

Deficiency in implicit reasoning has profound implications. It implies the models' limitations in inducing structured and compressed representations of facts and rules, which lead to redundant knowledge storage and difficulty in propagating knowledge updates~\citep{zhong-etal-2023-mquake}, and importantly, fundamentally impede the model from \textit{systematic generalization} over knowledge~\citep{lake2018generalization}. While explicit verbalizations of reasoning steps (e.g., chain-of-thought rationales) can improve task performance~\citep{wei2022chain,wang-etal-2022-iteratively,zelikman2022star,sun2023recitationaugmented,liu-etal-2023-crystal}, they are not available during large-scale (pre-)training where the model's core capabilities are acquired~\citep{zhou2023lima, lin2024the}.

\begin{figure}[ht]
\centering
\includegraphics[width=0.93\textwidth]{fig/fig1.pdf}

\caption{We find that transformers can learn to reason implicitly, but this skill is only robustly acquired through \textit{grokking}, i.e., an extended period of training far beyond overfitting. Moreover, the transformer fails to \textit{systematically} generalize for composition, yet succeeds for comparison. We conduct a mechanistic study into the model internals throughout grokking, which reveals distinct generalizing circuits across the two tasks (Figure~\ref{fig:circuit_composition}, \ref{fig:circuit_comparison}) that explains the variations in systematicity.}
% 
\label{fig:grokking}
\end{figure}

\textit{Is implicit reasoning doomed given that even the most capable models struggle? Can it be resolved by further scaling data and compute, or are there fundamental limitations of the transformer~\citep{vaswani2017attention} that prohibit robust acquisition of this skill?}

In this paper, we rigorously study these questions by constructing synthetic training and evaluation datasets, training transformers from scratch, and examining their generalization. We conceptualize reasoning as the \textit{induction and application of inference rules}, and expose the model to a mixture of ``atomic facts'' and ``inferred facts'' (which are deduced from the atomic facts via a set of latent rules), resembling ``axioms'' and ``theorems'' in a formal system. To evaluate how well the model learns the rules, we test its ability to make novel deductions (i.e., completing unseen inferred facts) in both in-distribution (ID) and out-of-distribution (OOD) scenarios.\footnote{Definitions of ID/OOD are introduced in \S\ref{sec:general_setup}.} This approach allows us to control the training data and perform clean evaluations, which would be challenging when studying existing LLMs trained on uncontrolled data.

Our experiments reveal that transformers \textit{can} learn to perform implicit reasoning, but this skill is only robustly acquired through \textit{extended training far beyond overfitting} (Figure~\ref{fig:grokking}), a phenomenon known as \textit{grokking}~\citep{power2022grokking}. We find that the speed of improvement in generalization correlates with the \textit{ratio} between inferred and atomic facts in training, and depends little on the absolute \textit{size} of the training data (Figure~\ref{fig:composition}). This suggests a correction of prior explanations of grokking based on \textit{critical data size}~\citep{liu2023omnigrok, varma2023explaining, zhu2024critical, huang2024unified}, in that it should instead be the \textit{critical data distribution} that decides the characteristics of grokking. Our findings extend prior observations of the grokking phenomenon primarily in algorithmic and linguistic tasks \citep{power2022grokking, murty-etal-2023-grokking} to the domain of knowledge-based reasoning, and deepen our understanding of the grokking phenomenon.

Moreover, we find that the transformer exhibits different levels of systematicity across reasoning types. While ID generalization is consistently observed, in the OOD setting, the model fails to systematically generalize for composition but succeeds in comparison (Figure~\ref{fig:grokking}). To understand why this happens, we conduct mechanistic analysis of the internal mechanisms of the model. The analysis uncovers the gradual formation of the \textit{generalizing circuit} throughout grokking and establishes the connection between systematicity and its configuration, specifically, the way atomic knowledge and rules are stored and applied within the circuit. Our findings imply that proper cross-layer memory-sharing mechanisms for transformers such as memory-augmentation~\citep{sukhbaatar2015end, graves2016hybrid} and explicit recurrence~\citep{dehghani2018universal, hutchins2022blockrecurrent, tan-etal-2023-sparse} are needed to further unlock transformer's generalization.

Finally, to demonstrate the power and potential of parametric memory for complex reasoning, we show that for a reasoning task with a large search space, a fully grokked transformer can achieve near-perfect accuracy, while state-of-the-art LLMs like GPT-4-Turbo~\citep{openai_devday_2023} and Gemini-1.5-Pro~\citep{google2024gemini} based on non-parametric memory fail badly regardless of prompting styles or retrieval augmentation.